% !TEX root = ../Grunddatei.tex
\section{Experimental results}

To execute our tool you just need to go into the
\texttt{OneWayStreet\_analysis\_complete}-folder and execute \texttt{run.py}.
Before that you should give the tool writing access to this folder so it can
create the results and save them. You also probably need to create your own
API Key at  for the access to \textit{MapTiler} and insert it as \texttt{api\_key}.
You can set the coordinates as \texttt{lat\_sat, lon\_sat} at the top of the script
as well as choose your own identifier for the results. Another parameter
important for testing is the radius. This variable tells our program how far to
search for streetview panoramas.
We recommend to choose your test coordinates with care in
\texttt{https://panoramax.openstreetmap.fr/} and set as a filter 360° images only.
It is generally best to actually pick coordinates where there is only one
intersection visible and where there are actually usable streetview images.

After starting the program it will first download the image with the exact
coordinates given as center and with given zoom and tile radius from
\textit{MapTiler}.
Then it uses the \textit{YOLO11} model to get the center point of the
intersection in the image and return its coordinates as latitude and longitude.

These values can then be used to download the nearest streetview panorama from
the given panoramax instance. Currently there is no API Key needed to
download images from the instance in the code, so we recommend just using it.
If this instance should however not be responsive there are some other instances,
whose urls can be used instead.

After this our chosen street detection function is used to crop the streetview
panorama for separate streets.

Then again for every crop the \texttt{classify\_street}-function is used to
classify the street as explained in methodology.

As in between results the satallite image, the segmented satellite image,
the streetview panorama and the crops of the streetview panorama are written
to the file system.

As a final result the \texttt{summary.json} is saved, which contains some metainformation
about the panoramax image, the coordinates from the start, the intersection center point
and finally the angles of the roads detected inside streetview panorama with
their respective label. This result is machine readable and could be resused by
other functions.

The following section is about the results of each subtask and how well they
perform as parts of the proposed workflow.
The first step is getting a satellite image at given coordinates.
This step works flawlessly and can be done easily, since it only requires to
define latitude and longitude. During our work in this project, we also had no
problems with the \textit{MapTiler} API, as we could stay below the free limit.
It must be mentioned, that the quality of satellite images, while very good for
the area around Caen, can vary depending on the given coordinates.
This is, for example, because available data originates from different
satellites, that could have worse quality.

The results of the second step, the intersection segmentation, are also quite good.
When looking at the model training, both precision and recall are equal to around 0.8.
This means, that the current model is good at finding an intersection in a given
satellite image as well as segmenting its area.
Another metric for model quality is the mean average precision, which equals to
around 0.5.
This means, that the model is slightly worse at creating a very detailed
segmentation of the intersection. In practice this means, that the resulting
segmentation might not reflect the actual area of the intersection accurately
and instead, creates segmentation areas that are fuzzy or include small parts of
adjacent areas. Fortunately, the effects on our proposed workflow are minimal,
since we don’t need or use a very detailed segmentation of a given intersection
and only use it to estimate the center point.

Third, this estimation of the intersection center point as well as the following
download of an adjacent street view panorama from \textit{Panoramax} also works
faultless.
However, it must be mentioned that the download can still fail, depending on the
chosen search radius and the density of street view panoramas.
When looking at different areas, it could be needed to change this search radius
parameter.
During our test for Caen, the current value did not disrupt the current implementation.
The next step of our workflow is the identification of relevant sections of a
given street view panorama.
The results show, that the proposed approach works in principle, but still needs
fine tuning.
During our tests, we only achieved a precision of 65.79 \% and a recall of
around 53 \%.
The F1-Score of our solution equals to 57.45 \%.
A reason for the current flaws is that the used street view panoramas can vary
immensely. Different shortcomings are, for example, a disparate recording angles
as well as the complexity of a given intersection scene, that can include
different buildings, structures, participants or other objects.
The current implementation does not account for all of those influences, which
is why irrelevant sections of the panorama might be identified incorrectly,
which impacts the continued workflow and the end results.

The fourth step is the detection of street signs in images that represent the
relevant sections of the streetview panorama. Our tests show, that the current
implementation produces accurate results and can be used as intended in our
proposed workflow. Here, the developers of the used models published own metrics
of their respective models on \textit{Hugging-Face}.
The \texttt{Panoramax/detect\_face\_plate\_sign}-model has a provided accuracy for
road signs of 0.879 which is still very high and the
\texttt{Panoramax/classify\_fr\_road\_signs}-model for classification has an
accuracy of 0.987. These results match our own tests during this project.

The last step of creating a results JSON file as well as all additional
(auxiliary) functions work flawlessly and produce the necessary data needed for
the following subtasks of the proposed workflow.


